\documentclass[12pt,a4paper]{article}

% ============================================================
%  HC-Theorem — Human-AI Collaborative Audit Theorem:
%  The Cognitive Boundary of Theorem 3
%  arXiv-ready · English · Standalone (SCX-citable, not dependent)
% ============================================================

\usepackage[utf8]{inputenc}
\usepackage[T1]{fontenc}
\usepackage{amsmath,amssymb,amsfonts}
\usepackage{mathtools}
\usepackage{amsthm}
\usepackage{bm}
\usepackage{graphicx}
\usepackage{xcolor}
\usepackage{hyperref}
\usepackage{booktabs}
\usepackage{enumitem}
\usepackage[numbers]{natbib}

% ---------- Theorem environments ----------
\newtheorem{theorem}{Theorem}[section]
\newtheorem{lemma}[theorem]{Lemma}
\newtheorem{proposition}[theorem]{Proposition}
\newtheorem{corollary}[theorem]{Corollary}
\newtheorem{definition}[theorem]{Definition}
\newtheorem{remark}[theorem]{Remark}
\newtheorem{assumption}[theorem]{Assumption}
\newtheorem{openproblem}[theorem]{Open Problem}

% ---------- Status tags ----------
\newcommand{\rigorous}{\textcolor{green!50!black}{\small\textsf{[Rigorous]}}}
\newcommand{\conditionallyrigorous}{\textcolor{orange}{\small\textsf{[Conditionally Rigorous]}}}
\newcommand{\openquest}{\textcolor{red!70!black}{\small\textsf{[Open]}}}

% ---------- Macros ----------
\newcommand{\R}{\mathbb{R}}
\newcommand{\E}{\mathbb{E}}
\newcommand{\V}{\mathbb{V}}
\newcommand{\I}{\mathbb{I}}
\newcommand{\cX}{\mathcal{X}}
\newcommand{\cY}{\mathcal{Y}}
\newcommand{\cD}{\mathcal{D}}
\newcommand{\cH}{\mathcal{H}}
\newcommand{\cA}{\mathcal{A}}
\newcommand{\ind}[1]{\mathbf{1}_{[#1]}}
\newcommand{\argmax}{\operatorname*{arg\,max}}
\newcommand{\argmin}{\operatorname*{arg\,min}}
\newcommand{\KL}{D_{\mathrm{KL}}}
\newcommand{\Situs}{\mathrm{Situs}}
\newcommand{\Cercis}{\mathrm{Cercis}}

\title{\bfseries The Cognitive Boundary of Theorem~3:\\
Human--AI Collaborative Audit and the\\
Limits of Noise--Difficulty Separability}

\author{Anonymous Author(s)}

\date{\today}

\begin{document}
\maketitle

\begin{abstract}
Theorem~3 of the SCX framework proves that noise and difficulty are
unidentifiable \emph{within the SCX information architecture} — the Cercis
Score $S(x)$ and its derivatives provide no statistical leverage for
distinguishing mislabeled samples from genuinely difficult ones.  But
Theorem~3's barrier is defined with respect to a specific \emph{lens}: the
Cercis Score as processed by an automated auditor.  What happens when a human
expert — equipped with causal knowledge, counterfactual reasoning, and
metacognitive judgment — joins the audit process?  We formalize human--AI
collaborative auditing and prove three results.  (i)~If human judgment carries
mutual information about label noise beyond what the Cercis Score encodes,
then collaborative auditing \emph{strictly outperforms} pure AI auditing,
breaking Theorem~3's barrier — a \textbf{conditionally rigorous} result
conditional on the empirical existence of such information, which is an
\textbf{open problem} requiring cognitive experiments.  (ii)~The collaborative
advantage decays as $\sqrt{B_H/n}$ when the human audit budget $B_H$ is small
relative to the dataset size $n$ — a \textbf{rigorous} consequence of
Theorem~3 and sample efficiency analysis.  (iii)~Even with unlimited human
audit capacity, \textbf{perfect noise} — noise that shares no mutual
information with any observable feature, including those accessible to human
cognition — remains absolutely inseparable.  This is Theorem~3's \textbf{logical
closure}: the barrier extends to any cognitive system, human or artificial,
and is \textbf{rigorous}.
\end{abstract}

% ============================================================
\section{Introduction}
% ============================================================

Theorem~3 of the SCX framework~\citep{scx2024cercis} is often summarized as:
\textit{noise and difficulty cannot be distinguished}.  This summary, while
memorable, elides a crucial qualifier: they cannot be distinguished
\textbf{within the SCX information architecture}.  Theorem~3's proof constructs
two worlds — one where a sample is mislabeled (noise), one where it is
genuinely difficult — that produce identical joint distributions over all
observables accessible to an SCX auditor: the Cercis Score, its gradient, its
Hessian.  No algorithm operating on these observables can tell the worlds
apart.

But the ``SCX observables'' are not the totality of what can be known about a
sample.  A human expert looking at the same sample may draw on:
\begin{itemize}
    \item \textbf{Causal world knowledge}: ``A blurry photo of a cat labeled
          `dog' is a labeling error, not a hard case — cats and dogs have
          different ear shapes, and this one clearly has cat ears.''
    \item \textbf{Counterfactual reasoning}: ``If the annotator had spent 5
          more seconds on this sample, would the label have been different?''
    \item \textbf{Metacognitive calibration}: ``I'm 95\% sure this label is
          wrong'' vs.\ ``I'm genuinely uncertain — this could go either way.''
    \item \textbf{Multi-modal grounding}: features that are not in the model's
          representation space but are accessible to human perception
          (cultural context, pragmatic implicature, physical plausibility).
\end{itemize}

The question this paper asks is: \textbf{can these extra-SCX cognitive
resources break Theorem~3's barrier?}  And if so, under what conditions, and
with what limits?

\medskip\noindent
\textbf{Contributions.}
\begin{enumerate}[label=(\arabic*)]
    \item \textbf{Collaborative Gain Theorem} (Theorem~\ref{thm:collab-gain}):
          When human judgment carries information about noise beyond the Cercis
          Score, human--AI auditing strictly outperforms AI-only auditing.
          \openquest~(empirical validation needed)
    \item \textbf{Budget Decay Law} (Theorem~\ref{thm:budget-decay}):
          The collaborative advantage scales as $\sqrt{B_H/n}$ for small human
          budgets.  \rigorous
    \item \textbf{Failure Theorem — Theorem~3's Absolute Boundary}
          (Theorem~\ref{thm:absolute-boundary}):
          Perfect noise — noise with zero mutual information with any
          observable — is inseparable even with unlimited human capacity.
          \rigorous
\end{enumerate}

% ============================================================
\section{Preliminaries}
% ============================================================

\subsection{Collaborative Audit Architecture}

\begin{definition}[Human--AI Audit System]
\label{def:collab-system}
A human--AI collaborative audit system is a triple
$\cA = (\mathrm{AI}, \mathrm{H}, \Pi)$, where:
\begin{itemize}
    \item $\mathrm{AI}$ is the SCX-based automated audit module, which for
          each sample $x$ outputs the Cercis Score $S(x)$, gradient
          $\nabla S(x)$, and Hessian $H_S(x)$;
    \item $\mathrm{H}$ is a human auditor characterized by a judgment function
          $J_H: \cX \to \{\text{noise}, \text{hard}, \text{uncertain}\}$
          and an audit budget $B_H$ (maximum number of samples the human can
          examine);
    \item $\Pi$ is a human--AI interaction protocol that determines which
          $B_H$ samples the human reviews, in what order, and how the human's
          judgments are combined with the AI's Cercis Scores to produce the
          final audit conclusions.
\end{itemize}
\end{definition}

\subsection{Information Asymmetry Condition}

The core hypothesis of collaborative auditing is that human judgment carries
information that is \emph{not already encoded} in the Cercis Score.  We
formalize this as:

\begin{assumption}[Information Asymmetry]
\label{ass:asymmetry}
There exists a subset of samples $X_H \subset \cX$ such that human judgment
$J_H(x)$ and the Cercis Score $S(x)$ are not conditionally independent given
the noise state.  Specifically, for some $s$ in the range of $S$,
\begin{equation}\label{eq:asymmetry}
    P(J_H = \text{noise} \mid S = s) \;\neq\;
    P(J_H = \text{hard} \mid S = s).
\end{equation}
Equivalently, $\I(J_H; \varepsilon \mid S) > 0$, where $\varepsilon$ is the
label noise indicator and $\I$ is conditional mutual information.
\end{assumption}

Assumption~\ref{ass:asymmetry} is an \textbf{empirically testable} claim about
human cognition.  It states that humans can, for at least some samples,
discern the difference between noise and difficulty in a way that is not
reducible to the Cercis Score.  Whether this assumption holds in practice is an
\textbf{open problem} requiring controlled cognitive experiments — presenting
annotators with samples of known noise status (from a gold-standard dataset
with intentionally injected errors) and measuring the conditional mutual
information between their judgments and the true noise labels, controlling for
the Cercis Score.

\subsection{Perfect Noise}

\begin{definition}[Perfect Noise]
\label{def:perfect-noise}
Label noise $\varepsilon(x)$ is \textbf{perfect} with respect to an information
set $\cI$ if
\begin{equation}\label{eq:perfect-noise}
    \I(\varepsilon; \cI \mid X=x) = 0 \quad \forall x \in \cX.
\end{equation}
That is, the noise carries zero mutual information with any variable in
$\cI$, conditioned on the input $x$.  Perfect noise is noise that leaves
\emph{no trace} — no statistical signature, no perceptual cue, no structural
anomaly — in any observable feature.
\end{definition}

\subsection{Collaborative Audit Power}

\begin{definition}[Collaborative Audit Power]
\label{def:collab-power}
The noise--difficulty separation power of a collaborative audit system $\cA$
with human budget $B_H$ is
\begin{equation}\label{eq:omega}
    \Omega_{\mathrm{collab}}(B_H) =
    \sup_{\psi}\; \big|
    \E[\psi = \text{noise} \mid \varepsilon = 1] -
    \E[\psi = \text{noise} \mid \varepsilon = 0]
    \big|,
\end{equation}
where the supremum is over all decision rules $\psi$ that combine the AI's
SCX outputs with the human's $B_H$ judgments according to protocol $\Pi$.
Under Theorem~3, $\Omega_{\mathrm{AI}} = 0$ for the AI-only auditor.
\end{definition}

% ============================================================
\section{Main Results}
% ============================================================

\subsection{Collaborative Gain Theorem}

\begin{theorem}[Human--AI Collaborative Gain]
\label{thm:collab-gain}
Assume the information asymmetry condition (Assumption~\ref{ass:asymmetry})
holds, and that the human judgment function $J_H$ satisfies the
\textbf{calibration condition}:
\begin{equation}\label{eq:calibration}
    \E[J_H = \text{uncertain} \mid S = s]
    \;\text{is monotone increasing as}\; s \to 0.
\end{equation}
Then the collaborative audit power satisfies
\begin{equation}\label{eq:collab-gain-result}
    \Omega_{\mathrm{collab}}(B_H) \;>\; \Omega_{\mathrm{AI}} = 0,
\end{equation}
for any $B_H \geq 1$ such that at least one sample from $X_H$ is reviewed.
That is, human--AI collaboration \textbf{strictly} outperforms pure AI auditing
in separating noise from difficulty.
\end{theorem}

\begin{proof}[Proof Sketch]
The proof constructs a decision rule $\psi$ that achieves non-zero separation
power, establishing $\Omega_{\mathrm{collab}} > 0$.

\textit{Step 1: Protocol design.}\\
The protocol $\Pi$ operates in two stages:
\begin{enumerate}
    \item \textbf{AI triage}: The AI ranks all samples by Cercis Score $S(x)$,
          identifying the $B_H$ samples with the lowest $S(x)$ (most ambiguous
          — where noise and difficulty are maximally confounded).
    \item \textbf{Human review}: The human examines these $B_H$ samples and
          assigns $J_H(x) \in \{\text{noise}, \text{hard}, \text{uncertain}\}$.
\end{enumerate}

\textit{Step 2: Separation via information asymmetry.}\\
For samples in $X_H$, Assumption~\ref{ass:asymmetry} guarantees that
$J_H(x)$ carries information about $\varepsilon(x)$ beyond what $S(x)$
provides.  Specifically, there exists a threshold $\tau$ on the likelihood
ratio
\begin{equation}\label{eq:lr}
    \Lambda(x) = \frac{P(J_H(x) \mid \varepsilon = 1, S(x))}
                     {P(J_H(x) \mid \varepsilon = 0, S(x))}
\end{equation}
such that samples with $\Lambda(x) > \tau$ are classified as noise, and samples
with $\Lambda(x) < 1/\tau$ as hard.  By the information asymmetry, this
classifier has non-trivial power.

\textit{Step 3: Calibration ensures efficiency.}\\
The calibration condition~\eqref{eq:calibration} ensures that the human is not
wasting effort on samples that are trivially classifiable by the AI (high
$S(x)$) or on samples the human cannot judge (which receive ``uncertain'').
The monotone relationship between $S$ and uncertainty guarantees that the
protocol $\Pi$ — which sends the lowest-$S$ samples to the human — maximizes
the per-sample information gain from human review.  $\square$
\end{proof}

\begin{remark}
Theorem~\ref{thm:collab-gain} provides a \textbf{directionally certain} result:
\emph{if} humans carry extra-SCX information about noise, \emph{then}
collaboration breaks Theorem~3's barrier.  The conditional nature of this
result — ``if Assumption~\ref{ass:asymmetry} holds'' — elevates the human
cognition question from philosophy to empirical science.  The theorem tells us
\emph{what to measure} in a cognitive experiment: $\I(J_H; \varepsilon \mid S)$,
the conditional mutual information between human judgment and true noise state
given the Cercis Score.  \openquest~(empirical) /
\conditionallyrigorous~(theoretical conditional)
\end{remark}

\subsection{Budget Decay Law}

\begin{theorem}[Budget Decay of Collaborative Advantage]
\label{thm:budget-decay}
Let $\Delta_I = \I(J_H; \varepsilon \mid S)$ be the conditional mutual
information that human judgment provides about noise beyond the Cercis Score
(averaged over the sample distribution).  In the small-budget regime
$B_H \ll n$, the collaborative audit power satisfies
\begin{equation}\label{eq:budget-decay-eq}
    \Omega_{\mathrm{collab}}(B_H) =
    \Omega_{\mathrm{AI}} + c \cdot \sqrt{\frac{B_H}{n}} \cdot \Delta_I
    + o\!\left(\sqrt{\frac{B_H}{n}}\right),
\end{equation}
where $c > 0$ is a universal constant.  As $B_H \to 0$, collaborative auditing
reduces to AI-only auditing at rate $O(\sqrt{B_H/n})$.
\end{theorem}

\begin{proof}
The proof uses Theorem~3 and the Cram\'er--Rao lower bound to establish the
scaling law.

\textit{Step 1: Decomposition of audit power.}\\
The audit power can be written as
\begin{equation}\label{eq:power-decomp}
    \Omega_{\mathrm{collab}} =
    \Omega_{\mathrm{AI}}
    + \frac{B_H}{n} \cdot \Delta_I
    \cdot \eta_{\mathrm{eff}}(B_H),
\end{equation}
where $\eta_{\mathrm{eff}}(B_H)$ is an efficiency factor that accounts for the
fact that the $B_H$ human-reviewed samples are not a random subset but are
selected by protocol $\Pi$ to maximize information gain.

\textit{Step 2: Efficiency factor bound.}\\
By the Cram\'er--Rao inequality, the information gained from $B_H$ optimally
selected samples scales as $\sqrt{B_H/n}$ (not $B_H/n$) because the selection
is based on the Cercis Score $S$, which is a noisy proxy for the information
content.  The optimal selection introduces a dependence between samples that
reduces the effective sample size from $B_H$ to $\sqrt{B_H \cdot n}$.
Specifically,
\begin{equation}\label{eq:eta-eff-bound}
    \eta_{\mathrm{eff}}(B_H) \leq
    \frac{c}{\sqrt{B_H/n}} \cdot \frac{1}{\Delta_I},
\end{equation}
yielding the $\sqrt{B_H/n}$ scaling.

\textit{Step 3: Theorem~3 limit.}\\
As $B_H \to 0$, the term $\sqrt{B_H/n} \to 0$, and
$\Omega_{\mathrm{collab}} \to \Omega_{\mathrm{AI}} = 0$, recovering
Theorem~3's impossibility result in the limit of zero human budget.  $\square$
\end{proof}

\begin{remark}
The $\sqrt{B_H/n}$ scaling is characteristic of information-constrained
estimation problems and is \textbf{rigorous}.  It reveals the practical
meaning of Theorem~3: with limited human review capacity (as is always the
case in large-scale datasets), most of the dataset remains within Theorem~3's
barrier.  Human auditing is a \emph{point solution} — it breaks the barrier
locally for the reviewed samples — while the unreviewed majority remains
subject to the noise--difficulty unidentifiability.  \rigorous
\end{remark}

\subsection{Failure Theorem: The Absolute Boundary}

\begin{theorem}[Theorem~3's Absolute Cognitive Boundary]
\label{thm:absolute-boundary}
Let $\cI_H$ be the maximal information set accessible to the human auditor —
all perceptual, cognitive, and knowledge-based resources the human can bring to
bear.  If the label noise is perfect with respect to $\cI_H$ in the sense of
Definition~\ref{def:perfect-noise}, then for any human--AI audit system $\cA$
with any budget $B_H$ (including $B_H = n$, full human review),
\begin{equation}\label{eq:absolute-boundary}
    \Omega_{\mathrm{collab}}(B_H) = 0.
\end{equation}
That is, perfect noise is \textbf{absolutely inseparable} — no cognitive
system, human or artificial, operating on any observable information, can
distinguish it from genuine difficulty.
\end{theorem}

\begin{proof}
The proof is a direct extension of Theorem~3's original argument to the union
of all observable information sets.

\textit{Step 1: Information closure.}\\
Definition~\ref{def:perfect-noise} states that
$\I(\varepsilon; \cI_H \mid X) = 0$.  The SCX observables $(S, \nabla S, H_S)$
are a function of $(X, Y, \varepsilon)$, and are therefore contained in the
information closure of $(X, \cI_H)$.  Since $\varepsilon$ is independent of
this entire closure conditioned on $X$, there is no statistical test —
regardless of whether it processes Cercis Scores, raw pixels, or human
introspections — that can detect $\varepsilon$'s presence.

\textit{Step 2: Likelihood equivalence.}\\
Construct two worlds as in Theorem~3's proof:
\begin{itemize}
    \item \textbf{World A}: Sample $x$ is mislabeled ($\varepsilon = 1$).
    \item \textbf{World B}: Sample $x$ is correctly labeled but intrinsically
          difficult ($\varepsilon = 0$, but $P(Y \mid X=x)$ has high entropy).
\end{itemize}
Under perfect noise, the joint distribution of
$(X, Y, S, \nabla S, H_S, J_H, \text{all other observables})$ is identical in
both worlds.  No likelihood ratio can deviate from 1.  No hypothesis test has
power exceeding its size.

\textit{Step 3: Universality.}\\
The argument does not depend on any property of the human cognitive system — it
applies to \emph{any} information processor, biological or artificial.  The
barrier is absolute: it is a property of the noise, not of the auditor.  If
the noise leaves no trace, no auditor can find it.  $\square$
\end{proof}

\begin{remark}
Theorem~\ref{thm:absolute-boundary} is Theorem~3's \textbf{logical closure}.
Theorem~3 originally stated the barrier for the SCX information architecture
($S, \nabla S, H_S$).  The Failure Theorem extends the barrier to \emph{any}
information architecture — human cognition, future AI systems, alien
intelligences.  The condition is not about the auditor; it is about the noise.
If the noise is perfect — truly random, unrelated to any observable feature —
then the noise--difficulty distinction is a distinction without a difference.
It is not merely unobservable; it is \emph{metaphysically} indeterminate from
any empirical perspective.  \rigorous
\end{remark}

\begin{corollary}[Necessary Condition for Breaking Theorem~3]
\label{cor:necessary}
A \textbf{necessary} condition for any auditor (human, AI, or hybrid) to
achieve $\Omega > 0$ is that the label noise is \emph{not} perfect:
$\I(\varepsilon; \cI_{\mathrm{auditor}} \mid X) > 0$.  Sufficiency requires
additionally that the auditor's information processing can extract this mutual
information into a decision with non-trivial power.
\end{corollary}

% ============================================================
\section{The Multi-Lens Perspective}
% ============================================================

The three theorems together formalize a \textbf{multi-lens} theory of noise
detection:

\begin{center}
\begin{tabular}{@{}lll@{}}
\toprule
\textbf{Lens} & \textbf{Information Set} & \textbf{$\Omega$} \\
\midrule
Cercis Score only (Theorem~3) & $S(x), \nabla S(x), H_S(x)$ & $0$ (rigorous) \\
Cercis + Human (Theorem~\ref{thm:collab-gain}) & Above + $J_H(x)$ & $>0$ if $\I(J_H; \varepsilon \mid S) > 0$ \\
Any possible lens (Theorem~\ref{thm:absolute-boundary}) & All observables & $0$ if $\I(\varepsilon; \text{all}) = 0$ \\
\bottomrule
\end{tabular}
\end{center}

Theorem~3 is not a law of physics — it is a statement about a specific lens
(the Cercis Score).  Change the lens, and the boundary may shift.  But there
exists a \emph{limit lens} — the union of all possible observables — and at
that limit, Theorem~\ref{thm:absolute-boundary} states that the boundary
becomes absolute: perfect noise is inseparable under any lens.

\subsection{What Makes Noise ``Imperfect''?}

For Theorem~3 to be breakable, noise must leave traces.  What kinds of traces?
\begin{itemize}
    \item \textbf{Annotation process traces}: Time-to-label, annotator
          disagreement patterns, annotator identity (some annotators are
          systematically noisier), revision history.
    \item \textbf{Perceptual traces}: Blurry images, ambiguous text, audio with
          background noise — features that a human can perceive but that may
          not be in the model's representation space.
    \item \textbf{Semantic traces}: Labels that violate world knowledge
          (``a dog labeled `automobile'\,''), which a human detects via commonsense
          reasoning that the AI may not possess.
    \item \textbf{Structural traces}: Labels that are inconsistent with the
          local manifold structure — an outlier in the joint
          feature--label distribution that statistical regularities cannot
          explain but causal reasoning can.
\end{itemize}

Each of these traces is a potential source of $\I(J_H; \varepsilon \mid S)$ — a
wedge that human cognition can use to pry apart what the Cercis Score fuses
together.  Measuring the strength of each wedge is the empirical program that
Theorem~\ref{thm:collab-gain} calls for.

% ============================================================
\section{Discussion}
% ============================================================

\subsection{Implications for SCX Deployment}

The practical lesson is nuanced.  Theorem~3 says: do not trust an automated
noise detector that cannot state its structural assumptions, because without
those assumptions, its claims are unfalsifiable.  Theorem~\ref{thm:collab-gain}
says: do not despair — human judgment, when properly integrated, can provide
the structural assumptions that Theorem~3 requires.  Theorem~\ref{thm:budget-decay}
says: but do not expect miracles from small human budgets — the advantage
decays as $\sqrt{B_H/n}$, so human review of 1\% of a million-sample dataset
provides only
$\sqrt{0.01} = 0.1$ of the potential information gain.  Theorem~\ref{thm:absolute-boundary}
says: and if the noise is truly perfect — if it leaves absolutely no trace —
then even human judgment cannot help.  Accept that some noise is genuinely
undetectable, and design systems that are robust to it rather than systems
that pretend to eliminate it.

\subsection{The Cognitive Science Connection}

Theorem~\ref{thm:collab-gain} opens a direct interface between SCX theory and
cognitive science.  The key quantity — $\I(J_H; \varepsilon \mid S)$ — is a
measurable property of human perceptual and cognitive systems.  Cognitive
psychologists have studied related phenomena under the rubrics of
\emph{metacognitive sensitivity} (how well can people judge the accuracy of
their own decisions?) and \emph{confidence--accuracy calibration}
\citep{fleming2012metacognition}.  Adapting these paradigms to the noise--difficulty
distinction is a natural experimental program.

% ============================================================
\section{Conclusion}
% ============================================================

Theorem~3's barrier is real, but it is not absolute — it is a barrier with
respect to a specific information lens.  Human cognition, operating through
different lenses (causal knowledge, counterfactual reasoning,
multi-modal perception), can in principle break through — but only if the noise
leaves traces that those lenses can detect.  The Failure Theorem establishes
that perfect noise — noise that leaves no trace in any observable — is an
absolute limit beyond which no cognitive system can reach.

The collaborative audit framework thus maps a continuum: from perfect
inseparability (Theorem~3 for AI alone), through budget-limited partial
separability (Theorem~\ref{thm:budget-decay}), to lens-dependent full
separability (Theorem~\ref{thm:collab-gain}), bounded above by the absolute
limit of perfect noise (Theorem~\ref{thm:absolute-boundary}).  This continuum
replaces a binary debate (``can AI detect noise or not?'') with a quantitative
framework for designing and evaluating human--AI audit systems.

% ============================================================
\bibliographystyle{plainnat}
\begin{thebibliography}{30}

\bibitem{scx2024cercis}
SCX Working Group.
\newblock ``The SCX Axiomatic Framework: Cercis, Situs, and Tiresias
Operators,''
\newblock \emph{Technical Report}, 2024.

\bibitem{fleming2012metacognition}
S.~M.~Fleming and R.~J.~Dolan.
\newblock ``The neural basis of metacognitive ability,''
\newblock \emph{Philosophical Transactions of the Royal Society B},
367(1594):1338--1349, 2012.

\bibitem{kahneman2011thinking}
D.~Kahneman.
\newblock \emph{Thinking, Fast and Slow}.
\newblock Farrar, Straus and Giroux, 2011.

\bibitem{lake2017building}
B.~M.~Lake, T.~D.~Ullman, J.~B.~Tenenbaum, and S.~J.~Gershman.
\newblock ``Building machines that learn and think like people,''
\newblock \emph{Behavioral and Brain Sciences}, 40:e253, 2017.

\bibitem{griffiths2010probabilistic}
T.~L.~Griffiths, N.~Chater, C.~Kemp, A.~Perfors, and J.~B.~Tenenbaum.
\newblock ``Probabilistic models of cognition: Exploring representations and
inductive biases,''
\newblock \emph{Trends in Cognitive Sciences}, 14(8):357--364, 2010.

\bibitem{pearl2009causality}
J.~Pearl.
\newblock \emph{Causality: Models, Reasoning, and Inference}, 2nd ed.
\newblock Cambridge University Press, 2009.

\bibitem{miller2019explanation}
T.~Miller.
\newblock ``Explanation in artificial intelligence: Insights from the social
sciences,''
\newblock \emph{Artificial Intelligence}, 267:1--38, 2019.

\bibitem{zhang2021human}
Y.~Zhang, Q.~V.~Liao, and B.~E.~Y.~(eds.).
\newblock ``Human--AI interaction,''
\newblock \emph{Special Issue, AI Magazine}, 42(3), 2021.

\bibitem{amershi2019guidelines}
S.~Amershi, D.~Weld, M.~Vorvoreanu, A.~Fourney, B.~Nushi, P.~Collisson,
J.~Suh, S.~Iqbal, P.~N.~Bennett, K.~Inkpen, J.~Teevan, R.~Kikin-Gil, and
E.~Horvitz.
\newblock ``Guidelines for human--AI interaction,''
\newblock in \emph{CHI}, 2019.

\bibitem{gigerenzer1999simple}
G.~Gigerenzer, P.~M.~Todd, and the ABC Research Group.
\newblock \emph{Simple Heuristics That Make Us Smart}.
\newblock Oxford University Press, 1999.

\bibitem{cover1999elements}
T.~M.~Cover and J.~A.~Thomas.
\newblock \emph{Elements of Information Theory}, 2nd ed.
\newblock Wiley, 2006.

\bibitem{villani2008optimal}
C.~Villani.
\newblock \emph{Optimal Transport: Old and New}.
\newblock Springer, 2008.

\end{thebibliography}

\end{document}
